\origpage{57}{GRAPH MINORS}
If $i=i'$ the result is true, and so we assume that $i\ne i'$. Define
\begin{equation*}
Z=S_{\nu i'-\nu+1}\cup Z_{\nu i'-\nu+2}\cup\cdots\cup Z_{\nu i'-1}\cup S_{\nu i'}
\end{equation*}
and
\begin{equation*}
S=S_{\nu i'-\nu+1}\cup S_{\nu i'-\nu+2}\cup\cdots\cup S_{\nu i'-1}\cup S_{\nu i'}.
\end{equation*}
By (4.3), $y\in Z$ and there is a path of $G$ from $y$ to some $s\in S$ within $Z$.
Now $Z_j\cap Z_{j'}\ne\varnothing$ only if $|j-j'|\le 1$, and so this path avoids
\begin{equation*}
S_{\nu i-\nu+1}\cup Z_{\nu i-\nu+2}\cup\cdots\cup Z_{\nu i-1}\cup S_{\nu i}
\end{equation*}
since $i\ne i'$. Hence by (4.3) it avoids $Y_i$. Thus $Y_i$ separates $y$ and $\rho(G)$ if and
only if it separates $s$ and $\rho(G)$. The result follows from (4.4).

For $1\le i\le T-1$, let $W_i$ be the set of all vertices $v$ of $G$ such that $Y_i$
separates $v$ and $\rho(G)$, and $Y_{i+1}$ does not separate $v$ and $\rho(G)$. Let $W_0$ be the
set of vertices not separated from $\rho(G)$ by $Y_1$, and let $W_T$ be the set of
vertices separated from $\rho(G)$ by $Y_T$. Then the sequence
\begin{equation*}
W_0\cup Y_1, W_1\cup Y_2, W_2\cup Y_3,..., W_{T-1}\cup Y_T, W_T
\end{equation*}
is a path-decomposition of $G$ by (4.5) and (1.7), and the intersection of any
consecutive pair of terms of this sequence has cardinality ${\le}\frac12\theta^2$. In order to
complete step 2, it suffices to show that $G\,|\,W_0,..., G\,|\,W_T$ each have path-width
${\le}p(\gamma-1, \mu)$, by (1.6) and (3.6)(ii). By (1.4), it suffices to show that
for $0\le i\le T$, every component of $G\,|\,W_i$ has path-width ${\le}p(\gamma-1, \mu)$.

Thus, take $0\le i\le T$, and let $C$ be a component of $G\,|\,W_i$. Assign a root
to $C$, arbitrarily. By our initial hypothesis it is sufficient to show that $C$ has
no rooted minor isomorphic to $P(\gamma-1, \mu)$. Suppose for a contradiction, that
it does; then $C$ has a subgraph $D$ which is isomorphic to a subdivision of
$P(\gamma-1, \mu)$.

Let
\begin{equation*}
A=\bigcup_{\nu(i+1)\le j\le n}S_j,\quad B=\bigcup_{1\le j\le\nu(i-1)}S_j.
\end{equation*}
\stmt
(4.6). \textit{$A, B$ and $W_i$ are disjoint.}
\stmt
For every vertex of $A$ is separated from $\rho(G)$ by $Y_i$ (if $i>0$) and by $Y_{i+1}$
(if $i<T$), from (4.4). Every vertex of $W_i$ is separated from $\rho(G)$ by $Y_i$ (if
$i>0$) but not by $Y_{i+1}$ (if $i<T$), by the definition of $W_i$. Every vertex of $B$ is
separated from $\rho(G)$ neither by $Y_i$ (if $i>0$) nor by $Y_{i+1}$ (if $i<T$), from (4.4).
The result follows.
